% UTF-8. Build with XeLaTeX or Tectonic. Replace sample content for real use.
\documentclass[UTF8,a4paper,11pt,fontset=fandol]{ctexart}
\usepackage[margin=2cm]{geometry}
\usepackage{amsmath,amssymb,booktabs,tabularx,longtable,array}
\usepackage{tikz}
\usepackage[hidelinks]{hyperref}
\usepackage{enumitem}
\setlist{nosep,leftmargin=2em}
\setlength{\parindent}{0pt}
\setlength{\parskip}{0.35em}
\newcommand{\conditions}[1]{\par\textbf{条件：}#1\par}
\newcommand{\source}[1]{{\footnotesize\textit{依据：#1}}\par}
\newcommand{\keypoint}[1]{\par\textbf{#1}\quad}
\newcolumntype{Y}{>{\raggedright\arraybackslash}X}
\title{课程精要：可编译模板}
\author{}
\date{}
\begin{document}
\pagestyle{plain}
\maketitle
\vspace{-2.5em}
\textbf{使用说明：}本文件用一阶线性方程演示知识块、条件、方法与示意图的排版。
实际精要应替换标题、来源和内容；不保留无用栏目。

\section{一阶线性模型（排版示例）}
\subsection*{掌握要求}
能够识别常系数一阶线性初值问题，确定稳态值与时间常数，并检验解是否满足初值和方程。

\subsection*{核心关系与条件}
考虑模型
\begin{equation}\label{eq:model}
 \tau\frac{\mathrm{d}y}{\mathrm{d}t}+y=b,
 \qquad y(0)=y_0 .
\end{equation}
\conditions{$\tau>0$、$b$ 为常数，$t\geq0$；$y$ 为状态量，$y_0$ 与 $b$ 与其同量纲，$\tau$ 的量纲为时间。}
其解为
\begin{equation}\label{eq:solution}
 y(t)=b+(y_0-b)\exp(-t/\tau).
\end{equation}
\keypoint{最短推导}
令 $z=y-b$，则 $\tau z'+z=0$，由初值得
$z(t)=(y_0-b)\exp(-t/\tau)$。
\source{模板自带的原创教学示例；真实任务替换为已核实的教材节号与页码。}

\subsection*{辨析与方法}
\begin{tabularx}{\linewidth}{@{}YY@{}}
\toprule
判断 & 依据与边界 \\
\midrule
时间常数越大，趋近稳态越慢 & 仅在此稳定一阶模型且比较同一归一化偏差时成立。\\
$y(\tau)=b$ & 一般不成立；偏差为初始偏差的 $\mathrm{e}^{-1}$。若 $y_0=b$，则始终在稳态。\\
\bottomrule
\end{tabularx}

\keypoint{题型方法}
先将方程化为式\eqref{eq:model}，识别 $\tau,b,y_0$，再用式\eqref{eq:solution}；
检查 $y(0)=y_0$、$t\to\infty$ 时 $y\to b$，以及代回后的残差为零。

\begin{center}
\begin{tikzpicture}[x=0.9cm,y=0.50cm]
 \draw[->] (0,0)--(6.2,0) node[right] {$t/\tau$};
 \draw[->] (0,0)--(0,3.7) node[above] {$(y-y_0)/(b-y_0)$};
 \draw[dashed] (0,3)--(6,3) node[right] {$1$};
 \draw[thick,domain=0:5.8,samples=60] plot (\x,{3*(1-exp(-\x))});
 \node[below] at (1,0) {$1$};
 \draw[dotted] (1,0)--(1,1.896)--(0,1.896);
 \node[font=\footnotesize] at (4,1.3) {式\eqref{eq:solution}，$b\ne y_0$};
\end{tikzpicture}
\end{center}
\end{document}
